\documentclass[10pt,a4paper]{article}
\usepackage[utf8]{inputenc}
\usepackage{amsmath}
\usepackage{amsfonts}
\usepackage{amssymb}
\usepackage{graphicx}
\usepackage[left=2cm,right=2cm,top=2cm,bottom=2cm]{geometry}
\usepackage{listings}
\usepackage{color}
\definecolor{munsell}{rgb}{0.0, 0.5, 0.69}
\definecolor{dkgreen}{rgb}{0,0.6,0}
\definecolor{gray}{rgb}{0.5,0.5,0.5}
\definecolor{mauve}{rgb}{0.58,0,0.82}
\definecolor{minas}{RGB}{0.244, 0.172, 0.36}
\definecolor{PUJ}{RGB}{44, 86, 151}
\definecolor{PUJ2}{RGB}{43, 93, 156}
\definecolor{PUJ3}{RGB}{20, 52, 107}
\definecolor{cyandk}{rgb}{0.0, 0.72, 0.92}

\author{Iván Andrés Trujillo }

%% code information
\lstset{frame=tb,
  language=Python,
  aboveskip=3mm,
  belowskip=3mm,
  showstringspaces=false,
  columns=flexible,
  basicstyle={\small\ttfamily},
  numbers=none,
  classoffset=1,
  morekeywords={True,False}, keywordstyle=\color{munsell}, 
  classoffset=0, 
  keywordstyle=\color{blue},  
  commentstyle=\color{dkgreen},
  stringstyle=\color{PUJ3},
  numberstyle=\tiny\color{gray},
  breaklines=true,
  breakatwhitespace=true,
  tabsize=4,
}

\begin{document}

\newcommand{\HRule}{\rule{\linewidth}{0.5mm}}

%----------------------------------------------------------------------------------------
%	LOGO SECTION
%----------------------------------------------------------------------------------------

\includegraphics[width = 4cm]{pujshield.eps}\\[0.5cm] 

\begin{center} 
\textbf{\LARGE
Introduction to programming in python for data analysis}\\[0.2cm]
\emph{\LARGE First assigment }\\[0.3cm] 
\emph{Iván A. Trujillo A.} \\
\textsc{\Large Pontificia Universidad Javeriana
}\\[0.2cm] 
\textsc{\large Facultad de ingeniería}\\[0.5cm] 
\HRule \\[0.4cm]
\end{center}

\section{Exercise}
the function sum() it is a built-in function, that allow us sum up the elements of a list or some objects.
\begin{lstlisting}
sum([1,2,3,4]
\end{lstlisting}
you must write a function to add a sequence of numbers.

The following code could be useful
\begin{lstlisting}
def counter(lista):
	for x in lista:
		statements
\end{lstlisting}

also we can uses the \textbf{*argum}
that allow us specify a function without a specified number of parameters. 
\begin{lstlisting}
def Function(*arguments):
	for x in arguments:
		statements
\end{lstlisting}



\section{Exercise}

Define a function for the binomial distribution and show as change the function with different values of $p$.
\begin{equation}
P(X=x) = \binom{k}{x}(1-p)^{k-x}x^{k}
\end{equation}

\section{Exercise }

Create a program to compute the following problem, \emph{How many total arrangements there are objects  take k in k from a population of N, for this case there is replacement and the order does not matter.}



\section{Exercise}

Order the following list \textbf{lista}
\begin{lstlisting}
lista = [[1,2,3], [4,5,6], [0,1,0]]
\end{lstlisting}
in the following \textbf{listaO}
\begin{lstlisting}
listaO=[[3,2,1],[6,5,4],[0,1,0]]
\end{lstlisting}

to resolve the problem you need used \textbf{while} or \textbf{for} loop.

\section{Considerations}

\begin{lstlisting}
import matplotlib.pyplot as plt # it is used to graph 
import numpy as np # to create array-structures
plt.plot(x,y)
\end{lstlisting}

In the combinatorial problems you need remember that, a in a combination there is not replacement and the order does not matter. In permutations there are not replacement but the order matter.

\end{document}